% Generated draft: requires mathematical and source review.
\documentclass[11pt]{article}
\usepackage[margin=25mm]{geometry}
\usepackage{fontspec}
\setmainfont{Latin Modern Roman}
\usepackage{amsmath,amssymb,mathtools,unicode-math}
\setmathfont{Latin Modern Math}
\usepackage{enumitem,needspace,xurl,hyperref}
\setlength{\parindent}{0pt}
\setlength{\parskip}{5pt}
\setlength{\emergencystretch}{3em}
\newcommand{\N}{\mathbb{N}}
\newcommand{\Z}{\mathbb{Z}}
\newcommand{\Q}{\mathbb{Q}}
\newcommand{\R}{\mathbb{R}}
\newcommand{\C}{\mathbb{C}}
\newcommand{\F}{\mathbb{F}}
\newcommand{\A}{\mathbb{A}}
\newcommand{\PP}{\mathbb{P}}
\newcommand{\E}{\mathbb{E}}
\newcommand{\eps}{\varepsilon}
\newcommand{\dd}{\,\mathrm{d}}
\begin{document}
\section*{UnsolvedMath ID: 30001743}
\textbf{Source problem code:} catalogue entry
\subsection{Definitions and mathematical statement}
The supplied record asks about extending thermostat-based reduced modeling from ordinary differential equations to grid-based discretizations of partial differential equations. It does not specify a particular equation, discretization, thermostat, or approximation criterion.

For schematic coordinate bookkeeping only, let \(I\) denote the index set of fine-grid unknowns, and write \(x(t)=(x_i(t))_{i\in I}\) for a trajectory of the full discretized system. If \(R\subseteq I\) indexes retained unknowns, their coordinate projection is
\[
P_Rx(t)=(x_i(t))_{i\in R}.
\]
Under the usual retained/unresolved partition, the unresolved indices would be \(I\setminus R\). This notation does not supply state spaces, equations of motion, or hypotheses absent from the source.

\textbf{Source question.} Can a thermostat-based reduced model for such a grid discretization use the unresolved fine-grid unknowns as its thermodynamic reservoir, while representing the reduced model by a subset of the fine-grid unknowns? Can the full trajectory then be approximated weakly by the resulting reduced dynamics?

\textbf{Unresolved qualifications.} The source wording describes the reduced model as a subset of ``those unknowns,'' leaving unclear whether this refers to all fine-grid unknowns or specifically to the unresolved unknowns. Thus the retained/unresolved partition above is only a possible interpretation. The source also does not define what makes a collection of unknowns a thermodynamic reservoir, how the reduced dynamics or thermostat is constructed, or how a reduced trajectory is compared with a full trajectory of different dimension. No meaning of weak approximation, limiting parameter, time interval, admissible class of equations, or universal versus existential scope is specified. Consequently, the supplied evidence determines a research question but does not determine a single fully quantified mathematical proposition.
\subsection{Short English statement}
The question is whether unresolved fine-grid PDE variables can serve as a thermodynamic reservoir for a thermostat-based reduced model, and whether its dynamics can weakly approximate the full trajectory. The reservoir, retained variables, dynamics, and meaning of weak approximation require clarification before a precise theorem or conjecture can be stated.
\subsection{Sources}
\begin{itemize}
\item Supplied local UnsolvedMath registry record, permanent ID \(30001743\), source problem code \texttt{catalogue entry}, titled \emph{Thermodynamic Reservoirs in Grid-Based Reduced Models}; collection: Oberwolfach Reports Open Problems. The record gives proposed year 2011 and names no proposer.
\item Supplied source citation: \emph{Geometric Numerical Integration} (2011), Oberwolfach Reports, pp.~825--900, DOI 10.4171/owr/2011/16. Source URL and supplied literature source labeled ``original OWR report'': \url{https://doi.org/10.4171/owr/2011/16}. No passage-specific page or named contributor is supplied.
\item The supplied original statement contains the reservoir/subset question. The supplied corrected statement additionally asks about weak approximation of the full trajectory. The record labels its statement \texttt{corrected\_verified} and reports that the official report confirms both questions in the extension from ordinary differential equations to grid-based PDE discretizations. These are supplied verification claims, not independent verification in this draft.
\item Literature assessment dated exactly \texttt{2026-08-21}: the supplied curation labels the problem open and describes the extracted sentence as an isolated model-reduction question lacking enough definitions to identify a theorem or resolved claim. Its listed source checked is the original OWR report linked above. This is supplied dated curation, not an independently verified literature assessment; time-sensitive status and cited claims require independent checking.
\end{itemize}
\end{document}
